\section{Summary of Contributions} \label{sec:61-conclusion-summary}

\begin{foldable}
  The three contributions trace a progression of increasingly tight constraints, each answered by a more principled form of active fidelity enforcement.
  What unifies them is not a shared architecture but a shared design stance: when data is scarce, access is restricted, or tokens are discrete, fidelity to the source distribution must be built in, not assumed.
\end{foldable}

\begin{foldable}
  \textbf{DivBFKD} (Chapter~\ref{ch:30-research01}) addresses the few-shot black-box \ac{KD} problem, where the student has access to only a small unlabeled image set and may query the teacher solely for top-1 predictions.
  The key observation is that diversity of the training distribution---not merely its size---determines whether WGAN-synthesised samples can provide useful gradient signal.
  DivBFKD enforces fidelity by introducing high-confidence images, those the teacher classifies with high certainty, as real-image exemplars in the WGAN discriminator loop on-the-fly.
  Adaptive per-class confidence thresholds prevent the teacher's class-wise accuracy bias from skewing which images are admitted, ensuring that the selected exemplars represent the teacher's decision boundaries rather than its easiest cases.
  An ablation study confirms that removing adaptive thresholds---reverting to a global cutoff---degrades performance, validating that class-wise calibration of the filtering criterion is load-bearing, not incidental.
  DivBFKD achieves state-of-the-art performance among few-shot black-box \ac{KD} methods across multiple image classification datasets and architectures.
\end{foldable}

\begin{foldable}
  \textbf{DIPKD} (Chapter~\ref{ch:40-research02}) addresses the extreme case of access restriction: no real data at all, and a teacher that returns only top-1 labels.
  The fidelity problem becomes entirely generative---the student must be trained on a synthetic distribution that covers the teacher's decision surface without any real-image anchor.
  DIPKD addresses this through a three-phase pipeline where each phase serves a distinct purpose.
  The Synthesis phase constructs hierarchical multi-scale image priors from uniform noise with geometric transforms and semantic cutmixing, producing structural diversity before any teacher query.
  The Contrast phase passes these priors through a primer student backbone and applies a contrastive objective that optimises the synthetic dataset directly, pushing candidates apart in embedding space to prevent mode collapse.
  The Distillation phase then recovers soft probabilistic labels from the primer student---approximating the dark knowledge the black-box teacher withholds---and uses them alongside hard teacher labels to train the final student.
  DIPKD achieves state-of-the-art performance across twelve benchmarks, with the largest gains in fine-grained and domain-specific settings where synthetic diversity is most constrained.
\end{foldable}

\begin{foldable}
  \textbf{TAKE} (Chapter~\ref{ch:50-research03}) shifts the problem domain from image model compression to text dataset distillation~(\ac{DD}), and with it the distributional fidelity problem shifts from generation to selection.
  The goal is to replace a large labeled text corpus with a compact surrogate that preserves downstream task performance.
  TAKE is the first text \ac{DD} method to simultaneously satisfy three desiderata: knowledge-based weighting of samples by their task-aligned contribution, a dataset-level distillation objective that aligns the selected corpus with the full source distribution, and human-readable output that is auditable and transferable across architectures.
  Trajectory-aware influence scoring computes importance weights from reciprocal gradient norms accumulated across the training trajectory, correcting the hard-sample bias that convergence-point scoring introduces by assigning inflated importance to noisy or boundary examples.
  An \ac{LLM} generation pool produces diverse textual candidates, and discrete \ac{OT} via the Sinkhorn algorithm selects the subset that globally minimises the entropically regularised transport cost between the candidate pool and the source distribution.
  The importance-weighted \ac{OT} objective combines task alignment and distributional coverage as complementary objectives: importance alone directs the budget but ignores global structure; coverage alone ensures spread but ignores contribution; together they close failure modes that neither addresses individually.
  TAKE matches or exceeds prior text \ac{DD} baselines across six language benchmarks at extreme compression, establishing the first end-to-end text \ac{DD} method with theoretical guarantees.
\end{foldable}

\begin{foldable}
  Across all three contributions, the fidelity-first design stance enables deployment in regimes that scale-dependent assumptions foreclose: edge and embedded devices for DivBFKD and DIPKD, and low-resource, privacy-constrained, and carbon-sensitive language model development for TAKE.
  Model compression reduces inference cost and latency; dataset distillation reduces training expenditure, data-governance risk, and cumulative carbon footprint across full development lifecycles.
  The practical impact scales with the tightness of the constraint: the more restricted the access regime, the larger the relative gain from enforcing fidelity actively rather than assuming it.
\end{foldable}
